% ============================================================
% VI. CONCLUSIONES
% ============================================================

\section{Conclusiones}

Se implementó un sistema completo de digitalización de señales sobre un microcontrolador de 32 bits, cubriendo muestreo periódico a $500$ Hz con $12$ bits de resolución, transmisión eficiente de datos ($115200$ baudios, justificados numéricamente frente al límite insuficiente de $9600$ bit/s), dos métodos de reconstrucción (DAC directo y PWM + filtro pasa-bajos), y digitalización de dos tipos de sensores por I2C: un BME280 (presión, temperatura, humedad) a $100$ Hz —límite propio del sensor, no $200$ Hz, ya que la guía original nombra una IMU en su lugar— y un encoder de motorreductor a $200$ Hz, además de análisis en MATLAB para las cuatro tablas de resultados solicitadas por la guía, incluyendo la sección opcional de aliasing.

El análisis teórico de las tablas de resultados —posible sin hardware físico para las columnas de frecuencia normalizada calculada— mostró que $75$ Hz es la única de las cuatro frecuencias de la sección de digitalización periódica que no produce un número entero de muestras por ciclo ($6.667$), anticipando una mayor fuga espectral y, por tanto, una mayor diferencia esperada entre frecuencia normalizada calculada y observada para esa fila específica frente a $50$, $100$ y $125$ Hz. De igual forma, el análisis teórico de la sección de aliasing predijo que $550$, $950$, $1050$ y $2050$ Hz deberían observarse todos con la misma frecuencia normalizada ($0.100$, equivalente a $50$ Hz aparentes), al igual que $575$/$1075$/$2075$ Hz ($0.150$) y $500$/$1000$/$2000$ Hz ($0.000$, indistinguibles de continua) — la firma característica del plegamiento espectral que define el aliasing, y una predicción verificable directamente una vez existan capturas reales.

Ningún resultado de esta práctica proviene de mediciones físicas: no se dispuso de tarjeta ESP32-WROOM, osciloscopio, generador de señales, BME280 ni encoder/motorreductor en el entorno de desarrollo de este informe. El firmware, el programa de captura eficiente, y los scripts de análisis en MATLAB están completos, documentados y listos para ejecutarse; el trabajo pendiente directo es flashear los cinco sketches en una tarjeta ESP32-WROOM real, capturar los datos con el hardware de laboratorio (osciloscopio, generador de señales, BME280 y motorreductor con encoder), y reemplazar cada valor marcado como \emph{pendiente} en las Tablas~\ref{tab:periodica}–\ref{tab:aliasing} por su medición real, verificando en el proceso las diferencias predichas teóricamente en esta sección de Discusión.
